\section{Conclusions}

The numerical and experimental analyses provide complementary assessments of canonical internal and separated flows.

For laminar pipe flow, the final refinement changes the selected observables by less than $1\%$. The CFD solution reproduces the developed Poiseuille profile with a relative $L_\infty$ error of $3.18\%$ and the analytical pressure gradient within $1.24\%$. For turbulent pipe flow at $Re_b=75{,}000$, the standard $k$--$\varepsilon$ configuration gives the lowest profile RMSE and a Darcy friction factor of 0.018984, only $0.082\%$ from the experimental reference.

The force transducer is linear over the investigated range, with $R^2=0.9999988$ and a force RMSE of 0.00572~N. Near-zero loads require the longest settling time and define the practical limit for interpreting mean lift. The cylinder experiments give $\overline{C_D}=1.499\pm0.076$ and $\overline{C_L}=-0.028\pm0.057$ as standard-uncertainty results. A local velocity probe gives $St=0.204$, while the lift spectrum gives $St=0.181$ at a second operating condition. Both analyses identify a coherent alternating wake.

The combined results establish a consistent sequence from numerical-resolution assessment and analytical verification to turbulence-model validation, sensor calibration, velocity-field measurement, spectral identification, and uncertainty-aware interpretation.
